\chapter{Conclusion générale}

Ce stage a porté sur la conception et le développement de \textbf{Secure Requirements
Studio}, un outil interne d'ingénierie des exigences destiné à cadrer le futur projet de
Plateforme de Gestion des Bénéficiaires de l'Axe Économie Sociale et Solidaire de la
Fondation OCP, en intégrant dès l'origine des exigences de cybersécurité.

Le travail réalisé a permis de mettre en place un entretien structuré — assisté par IA
ou en questionnaire classique hors ligne — capable de transformer les réponses des
parties prenantes en un modèle de projet structuré (Knowledge Graph), puis en exigences
traçables, en une analyse de sécurité proportionnée aux risques réels, et en un Cahier
des Charges concis accompagné d'une conception applicative et d'un diagramme de
conception MVP unique.

Au-delà de l'aspect technique, ce stage a été l'occasion de travailler un principe
directeur exigeant~: ne jamais présenter une supposition ou une inférence de l'IA comme
un fait acquis de la Fondation OCP, et distinguer systématiquement ce qui est confirmé,
supposé, ou à confirmer. Ce principe a structuré aussi bien la conception du moteur de
génération d'exigences que celle du moteur de sécurité.

\TODO{Compléter cette conclusion avec un bilan personnel du stage (apprentissages
techniques et humains, difficultés rencontrées et comment elles ont été surmontées) une
fois le stage arrivé à son terme.}

Les livrables produits par Secure Requirements Studio — Cahier des Charges, conception
applicative, diagramme MVP et exigences de cybersécurité — constituent la base sur
laquelle l'équipe de développement pourra s'appuyer pour concevoir, dans un second
temps, la Plateforme de Gestion des Bénéficiaires elle-même.
